\section{Conclusion}
Finite-catalog joint renewal design has an exact path structure and a nontrivial complexity boundary. The original constraints yield constructive selected-boundary resources, and history-level component packing extends that structure to heterogeneous terminal rewards. Exact priced paths turn the same history-level supports into polynomial node bounds. Common expected caps admit an exact two-command reduction, whereas heterogeneous caps and paid activation already give NP-completeness with unit linear reward, zero service costs, and a nonbinding command budget.

These results support complementary algorithms rather than one universal solver claim. Uniform resource gridding gives a polynomial additive guarantee with explicit numerical accuracy dependence. Grid-free price branching gives compact, independently checked original-instance intervals and finite exact termination, with a potentially exponential tree. The experiments expose both favorable larger-catalog cases and difficult reduction instances. Promise-dependent forced-command bounds extend safe exclusion to anchors but need not repay their cost on every model.

The mathematical model still uses one aggregate obligation, a common ordered catalog with prefix eligibility, additive opening charges and pre-draw service. The heterogeneous extension does not establish a multiple-resource or nonprefix result. All reported operational inputs are synthetic; no field calibration or operational tolerance value is inferred. The theory, constructive recovery, and reproducible comparisons supply the methodological contribution.

\section*{Code and Data Accessibility}
The accompanying \texttt{CODE\_AND\_DATA.zip} contains the current sources, rational inputs, frozen protocols, execution records, exact certificates, numerical comparisons, and reproduction commands, together with unchanged baseline dependencies. Historical computational records remain at their original revision paths. Complete older readers remain in the repository; the current root-reader predecessors are also in the archive. The companion supplies retained proofs and diagnostics. The README explains execution, verification requirements, and the preservation and build manifests.
